% ============================================================
%  §6 Συμπεράσματα και μελλοντική έρευνα
% ============================================================

\label{sec:conclusions}

Η παρούσα εργασία παρουσίασε ένα ολοκληρωμένο πλαίσιο για την ενσωμάτωση θερμικής
συμπεριφοράς ελαστικών σε προσομοιωτή γύρου, με έμφαση στην αριθμητική αποδοτικότητα και
τη φυσική ερμηνευσιμότητα. Η μεθοδολογία επέτρεψε, για πρώτη φορά στο πλαίσιο αυτό, τη
συνεκτίμηση θερμοκρασιακής εξέλιξης και δυναμικής οχήματος σε έναν αλγόριθμο quasi-static.

\subsection{Κύρια συμπεράσματα}

Τα κυριότερα ευρήματα της εργασίας συνοψίζονται ως εξής. Πρώτον, ο ενσωματωμένος
προσομοιωτής (slip estimation → thermal ODE → wear model) παράγει ποιοτικά ρεαλιστικά
αποτελέσματα σε \textit{[TODO: χρόνος εκτέλεσης]} -- αρκετά γρήγορα για παραμετρική ανάλυση
αλλά σημαντικά αργότερα από έναν κλασικό point-mass lap sim. Δεύτερον, η θερμοκρασία
ελαστικού παρουσιάζει σαφή τοπική συσχέτιση με τη δυναμική πίστας: αργές στροφές τείνουν
να εντείνουν τη φθορά ανά μέτρο, ακόμα και αν η απόλυτη διαφορά θερμοκρασίας είναι μικρή.
Τρίτον, η επιλογή αρχικής θερμοκρασίας (out-lap vs. hot-lap) επηρεάζει σημαντικά τα
αποτελέσματα στα πρώτα \textit{[TODO]} km, επιβεβαιώνοντας τη σημασία του warm-up
strategy.

% TODO: Συμπλήρωσε με τα πραγματικά ποσοτικά αποτελέσματα όταν τα έχεις.
% Αυτό εδώ είναι σκελετός για την αφήγηση.

\subsection{Περιορισμοί της τρέχουσας υλοποίησης}

Η υλοποίηση φέρει αρκετούς περιορισμούς που πρέπει να αναφερθούν. Το single-layer θερμικό
μοντέλο αγνοεί τη θερμοκρασία του εσωτερικού αέρα (inner liner), που λειτουργεί ως
θερμικός reservoir και επηρεάζει ιδιαίτερα μακρές ευθείες. Επιπλέον, οι συνθήκες
περιβάλλοντος ($T_{\rm amb}$, $T_{\rm track}$) είναι hardcoded αντί για ορίσματα εισόδου.
Τέλος, ο αλγόριθμος δεν μοντελοποιεί το lap-over-lap buildup: κάθε γύρος ξεκινά από
αρχικές συνθήκες που δίνονται εξωτερικά, χωρίς αλληλεξάρτηση διαδοχικών γύρων.

\subsection{Προτάσεις για μελλοντική έρευνα}

Η επέκταση της παρούσας εργασίας μπορεί να πάρει διάφορες κατευθύνσεις. Η πιο άμεση είναι
η πειραματική επαλήθευση: σύγκριση προσομοιωμένης θερμοκρασιακής εξέλιξης με δεδομένα
οδήγησης από πίστα. Αυτό απαιτεί αισθητήρες θερμοκρασίας ελαστικού (infrared pyrometers ή
τερμοζεύγη) και ακριβή τηλεμετρία.

Σε θεωρητικό επίπεδο, ενδιαφέρον παρουσιάζει η επέκταση σε μοντέλο δύο ή τριών στρωμάτων
(carcass + tread + inner liner), που θα αυξήσει την πιστότητα σε βάρος υπολογιστικού
κόστους. Παράλληλα, ο εμπλουτισμός του μοντέλου φθοράς με δεδομένα από optical wear
sensors θα επέτρεπε in-loop βαθμονόμηση. Τέλος, η σύζευξη με αλγόριθμο real-time strategy
optimization (π.χ. MPC) θα αξιοποιούσε πλήρως τη δυνατότητα πρόβλεψης που παρέχει ο
προσομοιωτής.
